\subsection{Fase 6: Supervisión y mantenimiento}
En la literatura de MLOps \cite{kreuzberger2023mlops}, el monitoreo continuo forma parte del ciclo de operación de un modelo, mientras que en este trabajo se plantea solo como trabajo propuesto, por lo que se propone verificar los siguientes aspectos:

\begin{itemize}
    \item Calidad de las recomendaciones: Revisar si las peliculas recomendadas por el bot correspondan con lo pedido por el usuario.
    \item Comparacion entre metodos: Analisar las diferencias que hay entre los resultados mediante TF-IDF y embeddings, sobretodo cuando se hacen consultas en español.
    \item Consultas fuera del dominio: Verificar como es el comportamiento del modelo cuando recibe consultas que no esten relacionadas con peliculas
    \item Tiempo de respuesta: Medir cuanto tiempo le toma al sistema procesar la consulta y procurar que dicho tiempo no se exceda demasiado, ya que esto podria afectar la experiencia del usuario.
\end{itemize}

Para supervisar el funcionamiento del chatbot en Telegram, se propone verificar que los componentes del sistema funcionen correctamente. Esta supervision incluye: 

\begin{itemize}
    \item Recepcion de consultas, para verificar que el bot reciba correctamente los mensajes enviados por los usuarios.
    \item Funcionamiento de los metodos de recuperacion: Comprobar que los comandos \textit{/start}, \textit{/ayuda}, \textit{/metodo} y \textit{/asistente} ejecuten correctamente.
    \item Seleccion del metodo de recuperacion: Verificar que el bot utilice TF-IDF, embeddings o ambos metodos dependiendo la elegida por el usuario.
    \item Generacion de recomendaciones: Comprobar que las consultas sean procesadas de manera correcta y que se le presenten al usuario las peliculas recomendadas.
    \item Generacion de respuestas con LLM: Verificar que el comando \textit{/asistente} utilice los resultados de embeddings para generar una respuesta redactada en lenguaje natural.  
\end{itemize}